\chapter{Data Types}
\label{ch:types}

\begin{goals}
\begin{itemize}
  \item The four basic data types: integer, floating-point number, string and boolean
  \item How Salam works out the type of each value
  \item Integer division and how it differs from floating-point division
  \item Converting between types with \code{as}
  \item The whole family of number types, and when to choose each one
  \item Joining text and numbers together
\end{itemize}
\end{goals}

\section{Every value has a type}

In the last chapter we saw that the box of a variable can hold a number or a
piece of text. The computer, however, does not store numbers and text in the
same way, nor does it treat them the same way. Numbers can be added and
multiplied, text cannot; text can be read letter by letter, a number cannot.
For this reason every value in Salam has a \textbf{type}, and Salam always
knows what type of data is in each box.

In this book we deal with four basic types:

\begin{center}
\begin{tabular}{lll}
\toprule
Type & Name in Salam & Examples \\
\midrule
Integer & \code{int} & \code{17}, \code{-4}, \code{0} \\
Floating-point number & \code{f64} & \code{2.5}, \code{-0.75}, \code{3.0} \\
String (text) & \code{str} & \code{"hello"}, \code{"17"}, \code{""} \\
Boolean & \code{bool} & \code{true}, \code{false} \\
\bottomrule
\end{tabular}
\end{center}

These four are the types you will use most of the time, and \code{int} is
simply the short name of \code{i32}, a 32-bit integer. Salam has more
number types as well, and section \ref{sec:numtypes} introduces all of
them.

Salam works out the type of a value from the way you write it. You do not
have to write the type yourself:

\program{The four data types}
\begin{salamcode}
func main:
    count := 12          // an integer
    price := 2.5         // a floating-point number
    name := "book"       // a string
    available := true    // a boolean
    println count, price, name, available
end
\end{salamcode}

\begin{salamout}
12 2.5 book true
\end{salamout}

The rule is simple: a number without a decimal point is an integer; a number
with a point is floating-point; anything between quotation marks is a
string; and \code{true} and \code{false} are booleans. (The name
``boolean'' honours George Boole, the mathematician who first studied the
logic of true and false.)

\section{Integers}

An integer is a number with no fractional part: the number of students, a
page number, a year. It may be negative. The sum, difference and product of
two integers are always integers. Division, however, is another story:

\program{Dividing integers}
\begin{salamcode}
func main:
    println 7 / 2
    println 9 / 3
    println 1 / 4
end
\end{salamcode}

\begin{salamout}
3
3
0
\end{salamout}

Seven divided by two is three and a half, yet Salam printed \code{3}. Why?
Because both numbers were integers, and \textbf{dividing two integers in
Salam always gives an integer}: the fractional part is thrown away. That is
also why one divided by four came out as zero. This behaviour looks odd at
first, but it is extremely useful. If you share 7 apples between 2 people,
each person gets 3 whole apples and 1 apple is left over. Salam has a sign
for this ``left over'' as well: \code{\%}.

\program{Quotient and remainder}
\begin{salamcode}
func main:
    apples := 7
    people := 2
    println "Each person gets:", apples / people
    println "Left over:", apples % people
end
\end{salamcode}

\begin{salamout}
Each person gets: 3
Left over: 1
\end{salamout}

The remainder operator will serve us many times in the chapters ahead, for
instance to tell whether a number is even or odd.

\begin{note}[How large can an integer be?]
Every integer in Salam is kept in a box of a fixed size, and so it has a
limit. An \code{int} runs from about minus two billion to plus two
billion. For most programs this is more than enough, and for the rest
Salam has larger types, which section \ref{sec:numtypes} introduces.
\end{note}

\section{Floating-point numbers}

When we need numbers with a fractional part (a price, a weight, an average,
a temperature) we use floating-point numbers. All it takes is a decimal point
in the number. Even \code{3.0} is a floating-point number, although its value
is three:

\program{Arithmetic with floating-point numbers}
\begin{salamcode}
func main:
    println 7.0 / 2.0
    println 2.5 * 4.0
    println 10.0 / 4.0
    println 1.0 / 3.0
end
\end{salamcode}

\begin{salamout}
3.5
10
2.5
0.3333333333333333
\end{salamout}

There are two things to notice in this output. First, \code{2.5 * 4.0} is
\code{10.0}, but Salam printed \code{10}: when it prints a floating-point
number, Salam leaves out useless zeros at the end. Second, it printed one
third with sixteen decimal places, because one third has an endless decimal
expansion and the computer keeps only so much of it.

\begin{warn}
Floating-point numbers are stored in the computer \emph{approximately}. Run
this program:
\begin{salamcode}
func main:
    println 0.1 + 0.2
end
\end{salamcode}
\begin{salamout}
0.30000000000000004
\end{salamout}

This is not a fault in Salam; every computer in the world says the same,
because the number 0.1 has an endless expansion in base two (which is what
the computer works in), just as one third does in base ten. For everyday
work this tiny error does not matter at all, but remember that you should
never compare two floating-point numbers for being ``exactly equal''.
\end{warn}

\section{When integers and floating-point numbers mix}

If you add or divide an integer and a floating-point number, Salam converts
the integer to floating-point by itself and the result is floating-point:
\code{7 / 2.0} is \code{3.5}. But when \emph{both} sides are integers, the
division is integer division. If you want a floating-point result in that
case, you must convert at least one of the two sides to floating-point. For
this there is the word \kwd{as}, which converts a value to another type:

\program{Converting an integer to floating-point}
\begin{salamcode}
func main:
    total := 17
    count := 4
    average := (total as f64) / (count as f64)
    println "Average:", average
end
\end{salamcode}

\begin{salamout}
Average: 4.25
\end{salamout}

Without the conversion, \code{total / count} would have been a division of
two integers and would have given \code{4}. The parentheses are not
required, but they make it clear which value is being converted and make the
expression easier to read.

Conversion works in the opposite direction too. When you convert a
floating-point number to an integer with \code{as int}, its fractional part
is thrown away (it is cut off, not rounded):

\program{Converting floating-point to an integer}
\begin{salamcode}
func main:
    exact price := 3.75
    whole price := exact price as int
    println whole price
    negative := -3.75
    println negative as int
end
\end{salamcode}

\begin{salamout}
3
-3
\end{salamout}

\begin{note}
\code{as} works on any value, whether it is stored in a variable or written
right there: \code{3.75 as int} also gives \code{3}.
\end{note}

\section{Strings}

A string is text: a sequence of letters, digits and symbols written between
quotation marks. A string can be empty (\code{""}) and it can contain
digits, but a digit inside a string is not a number:

\program{String or number?}
\begin{salamcode}
func main:
    println 12 + 3
    println "12" + 3
    println "12" + "3"
end
\end{salamcode}

\begin{salamout}
15
123
123
\end{salamout}

Look at the second and third lines of the output. For strings, \code{+}
means ``join'', not ``add''. \code{"12" + 3} means: attach the number 3 to
the string ``12'', which gives the string ``123''. If either side of
\code{+} is a string, Salam converts the other side to a string as well and
joins the two. This feature is very useful for building messages:

\program{Building a message}
\begin{salamcode}
func main:
    name := "Sara"
    age := 17
    message := name + " is " + age + " years old."
    println message
end
\end{salamcode}

\begin{salamout}
Sara is 17 years old.
\end{salamout}

Mind the spaces: joining does not add any space of its own, so wherever we
wanted a space we included one inside the strings (\code{" is "}). Compare
this with \code{println name, age}, which puts a space between its items by
itself. Both ways are correct; with \code{+} you have more control over the
shape of the message, and with the comma you have less to type.

If a string holds digits and you want to do arithmetic with them, you must
convert it to a number. For this, strings have a capability called
\code{to\_int} that is attached to them with a dot:

\program{Converting a string to a number}
\begin{salamcode}
func main:
    text := "12"
    number := text.to_int()
    println number + 3
end
\end{salamcode}

\begin{salamout}
15
\end{salamout}

In the same way, \code{to\_float} converts a string to a floating-point
number. Chapter \ref{ch:strings} introduces the other capabilities of
strings.

\section{Booleans}

The boolean type has just two values: \code{true} and \code{false}. It may
look small, but every decision a program makes (chapter \ref{ch:conditions})
rests on these two values. Most of the time we do not write a boolean value
directly; it comes out of a comparison:

\program{Boolean values}
\begin{salamcode}
func main:
    age := 20
    is adult := age >= 18
    println is adult
    println age == 18
    println "hello" == "hello"
end
\end{salamcode}

\begin{salamout}
true
false
true
\end{salamout}

The sign \code{>=} means ``greater than or equal to'' and the sign
\code{==} (two equals signs in a row) means ``is it equal to?''. Note that
\code{==} is different from \code{=}: the first one asks, the second one
assigns. We will see all the comparison signs in the next chapter.

\section{The type of a variable does not change}

Once a variable is made with one type, it keeps that type to the end. You
cannot put text into a box that was made for a number:

\begin{salamcode}
func main:
    mut value := 10
    value = "ten"
    println value
end
\end{salamcode}

\begin{salamout}
error[E002]: cannot assign 'str' to 'i32'
 --> main.salam:3:5
\end{salamout}

The message says that a value of type \code{str} cannot be put into a
variable of type \code{i32} (the full name of \code{int}). Like every
other strictness of Salam, this one works in your favour: when you see a
variable called \code{count}, you can be sure that there is always a number
in it.

\section{The whole family of number types}
\label{sec:numtypes}

So far we have used two number types: \code{int} for whole numbers and
\code{f64} for numbers with a fractional part. They are the right choice for
most programs, and they are what Salam picks when you simply write
\code{17} or \code{2.5}. But Salam has a whole family of number types,
each with its own size. A larger box holds larger numbers but takes more
room in memory. The table below lists all of them.

{\small
\begin{longtable}{l>{\raggedright\arraybackslash}p{4.6cm}l}
\toprule
Type & What it holds & Range \\
\midrule
\endhead
\code{i8} & small whole number & $-128$ to $127$ \\
\code{i16} & whole number & $-32{,}768$ to $32{,}767$ \\
\code{i32}, \code{int} & whole number (the usual choice) & about $\pm 2.1$ billion \\
\code{i64} & large whole number & about $\pm 9.2 \times 10^{18}$ \\
\code{u8} & small whole number, never negative & $0$ to $255$ \\
\code{u16} & whole number, never negative & $0$ to $65{,}535$ \\
\code{u32}, \code{uint} & whole number, never negative & $0$ to about $4.3$ billion \\
\code{u64} & large whole number, never negative & $0$ to about $1.8 \times 10^{19}$ \\
\code{size} & a memory size that may be negative & the same as \code{i64} \\
\code{usize} & a size, never negative & the same as \code{u64} \\
\code{f32}, \code{float} & number with a fraction, about 7 digits & up to about $3.4 \times 10^{38}$ \\
\code{f64} & number with a fraction, about 16 digits & up to about $1.8 \times 10^{308}$ \\
\bottomrule
\end{longtable}}

The letter tells you the kind of number: \code{i} for an integer that may be
negative, \code{u} for an \emph{unsigned} integer, which is never negative,
and \code{f} for a floating-point number. The figure after it is the size
of the box in bits. Some types have a second, shorter name: \code{int} is
the same as \code{i32}, \code{uint} the same as \code{u32}, and
\code{float} the same as \code{f32}. So take care: Salam's everyday
floating-point type is \code{f64}, not \code{float}.

To give a value a particular type, write \code{as} and the type after it,
exactly as you did for conversions:

\program{Choosing a type with as}
\begin{salamcode}
func main:
    age := 17 as u8
    temperature := -12 as i8
    distance := 384400 as u32
    population := 8100000000
    price := 2.5 as f32
    println age, temperature, distance
    println population, price
end
\end{salamcode}

\begin{salamout}
17 -12 384400
8100000000 2.5
\end{salamout}

An age is never negative and never above 255, so \code{u8} suits it. A
winter temperature can be negative but stays small, so \code{i8} is enough.
The distance to the Moon in kilometres is large and never negative, so we
chose \code{u32}. Look at \code{population}: we gave it no type at all, but
the number is too large for an \code{int}, so Salam made it an \code{i64}
by itself.

\subsection{Going past the limit}

What happens if a calculation goes past the limit of its type? The number
does not grow; it wraps around and starts again from the other end, like
the odometer of an old car:

\program{Going past the limit}
\begin{salamcode}
func main:
    small := 250 as u8
    println small + 10
    big := 2147483647
    println big + 1
end
\end{salamcode}

\begin{salamout}
4
-2147483648
\end{salamout}

A \code{u8} cannot hold 260, so the result wrapped around to 4. In the same
way, one more than the largest \code{int} became the most negative
\code{int}. Salam gives no error for this, so choose a type that is large
enough for every value your program may produce. When in doubt, use
\code{int}, or \code{i64} for very large numbers.

\subsection{f32 or f64?}

\code{f64} keeps about 16 digits of a number and \code{f32} only about 7.
The difference shows when we store one third in both:

\program{f32 and f64}
\begin{salamcode}
func main:
    third := 1.0 / 3.0
    rough := third as f32
    println "f64:", third
    println "f32:", rough
end
\end{salamcode}

\begin{salamout}
f64: 0.3333333333333333
f32: 0.3333333432674408
\end{salamout}

The \code{f32} value is correct only up to its seventh digit; what comes
after that is noise. An \code{f32} takes half the memory of an \code{f64},
which matters when a program keeps millions of numbers, but for everyday
calculations \code{f64} is the better choice, and it is the type Salam
picks for \code{2.5}.

\subsection{Automatic conversions}

Salam converts a value from one number type to another by itself only when
nothing can be lost: an \code{int} into an \code{i64}, an \code{f32} into
an \code{f64}, or an \code{int} into an \code{f64}. In every other case you
must convert with \code{as}, even when the value would fit:

\begin{salamcode}
func main:
    steps := 120 as u8
    mut total := 0
    total = steps
    println total
end
\end{salamcode}

\begin{salamout}
error[E002]: cannot assign 'u8' to 'i32'
 --> main.salam:4:5
\end{salamout}

Writing \code{total = steps as int} fixes it. Salam asks you to be explicit
because a conversion between signed and unsigned types, or into a smaller
type, can change the value, as you saw above.

\begin{note}
There are a few more types that are not numbers. You have already met
\code{str} and \code{bool}. Chapter \ref{ch:strings} introduces
\code{char} for a single one-byte character and \code{uchar} for a single
Unicode character, and chapter \ref{ch:functions} introduces \code{void},
the type of a function that gives nothing back.
\end{note}

\section{A more complete program}

The weighted average of three courses with different credit weights, which
brings all the types of this chapter together:

\program{A weighted average}
\begin{salamcode}
func main:
    // grades are floating-point, credits are integers
    maths grade := 17.5
    maths credits := 3
    physics grade := 15.0
    physics credits := 2
    literature grade := 18.25
    literature credits := 2

    total credits := maths credits + physics credits +
        literature credits
    weighted total := (maths grade * (maths credits as f64)) +
        (physics grade * (physics credits as f64)) +
        (literature grade * (literature credits as f64))
    average := weighted total / (total credits as f64)

    println "Total credits:", total credits
    println "Average:", average
    println "Passed?", average >= 12.0
end
\end{salamcode}

\begin{salamout}
Total credits: 7
Average: 17
Passed? true
\end{salamout}

Two things to notice. First, we split a long calculation over several lines;
when a line ends with an operator such as \code{+}, Salam understands that
the expression is not finished and continues it on the next line. Second,
the average really is 17.0 (a weighted total of 119 divided by 7), and Salam
printed it as \code{17}, as it always does.

\begin{summary}
\begin{itemize}
  \item Every value has a type: \code{int} (a whole number), \code{f64} (a
        number with a fractional part), \code{str} (text) and \code{bool}
        (true or false). Salam tells the type from the shape of the value.
  \item Dividing two integers throws away the fractional part; if one side
        is floating-point, the result is floating-point. \code{\%} gives
        the remainder of a division.
  \item Floating-point numbers are approximate; do not compare them for
        exact equality.
  \item \code{value as type} converts a value to another type. Converting
        floating-point to an integer cuts off the fraction.
  \item \code{+} between a string and anything else means joining.
        \code{text.to\_int()} converts a string to a number.
  \item The type of a variable does not change after it is made.
  \item Salam has a whole family of number types: \code{i8} to \code{i64},
        \code{u8} to \code{u64}, \code{size}, \code{usize}, \code{f32} and
        \code{f64}. \code{int}, \code{uint} and \code{float} are short
        names for \code{i32}, \code{u32} and \code{f32}. Choose one with
        \code{value as type}, and pick a size large enough for every value.
\end{itemize}
\end{summary}

\section*{Exercises}
\markright{Exercises}

\exercise Without running them, say what each of these instructions prints,
then check: \code{println 10 / 3}, \code{println 10.0 / 3.0},
\code{println 10 \% 3}, \code{println "10" + 3}, \code{println 10 + 3}.

\exercise Put an air temperature in Celsius into a variable, convert it to
Fahrenheit and print it. The formula: Fahrenheit equals Celsius times 9
divided by 5, plus 32. Write the program once with a floating-point number
and once with an integer, and explain the difference between the two
outputs.

\exercise Put the total length of a film in minutes (say 135) into a
variable and, using division and remainder, print it in the form ``2 hours
and 15 minutes''.

\exercise Make three variables with the values \code{"3"}, \code{3} and
\code{3.0}. Print them, then add the number 2 to each one and print the
results. Guess the output before you run the program.

\exercise Put the price of an item in dollars (a floating-point number)
and the exchange rate from dollars to your own currency (an integer) into
two variables, and print the price of the item in your currency as a whole
number.

\exercise Choose a suitable type for each of these values, create a variable
for each with \code{as}, and print them all: the number of days in a year,
the number of seconds in a year, the population of a large country, the
weight of a parcel in kilograms, and the score of a game that is never
negative and never above 200.
